\documentclass[11pt]{article}
\usepackage{fontspec}
\usepackage[french]{babel}
\usepackage{amsmath}
\usepackage{amssymb}
\usepackage{amsthm}
\usepackage{mathrsfs}
\usepackage{tikz-cd}
\usepackage{enumitem}
\usetikzlibrary{babel}

\newcommand{\printedpage}[1]{\marginpar{\footnotesize p.~#1}}
\newcommand{\msp}[1]{\marginpar{\footnotesize [#1]}}

\begin{document}

\noindent\printedpage{54}\msp{page 108}
% head: A. Grothendieck : Dérivateurs, Ch. XVII, Édition des Universités de Montpellier II et Paris VII
\[
\begin{tikzcd}
X(1) \arrow[r, "X(\partial)"] \arrow[d, "u(1)"'] \arrow[dr, phantom, "\text{cocart.}"] & X(0) \arrow[d] \arrow[dd, dotted, bend left=40, "u(0)"] \\
Y(1) \arrow[r] \arrow[dr, dotted, "Y(\partial)"'] & \tilde{Y} \arrow[d, "\tilde{u}"'] \\
 & Y(0)
\end{tikzcd}
\qquad .
\tag*{(6.2.5.4)}
\]

Pour qu'on ait $u \in \Phi(\boldsymbol{\Delta}^1)$, il faut et il suffit que l'on ait
\[
u(0),\ u(1),\ \tilde{u} \ \in\ \Phi\ .
\tag*{(6.2.5.5)}
\]

J'ai déjà vérifié la chose (sous forme duale), mais ne retrouve plus l'endroit (peut-être n'ai-je pas pris des notes). Je n'ai pas envie de rechercher la démonstration ici.

Pour achever de prouver 6.1. d), i.e. l'égalité $C(I) = C_I$ pour $I$ artinien, j'ai envie de prouver encore quelque chose de plus générale que voici : % [sic] plus générale -> plus général, agreement with « quelque chose »

\textbf{Corollaire 6.3.} \textit{Soient $(\mathcal{M}, \Phi, \Psi)$ et $I \longmapsto \Phi(I) \subset \mathrm{Fl}(\mathcal{M}(I))$ comme dans \textup{6.2.5.} Soit $I$ un ensemble ordonné, et considérons aussi $\Phi_I$ défini (dans la situation duale) page \textup{25 (2.2)} (mais avec des notations duales \dots). Je dis que l'on a}
\[
\Phi_I \ \supset\ \Phi(I)\ ,
\tag*{(6.3.1)}
\]
\msp{page 109}\textit{et on a égalité si $I$ est noethérien, i.e. $I^{\circ}$ est artinien.}

\textsc{Démonstration}. La démonstration de $\Phi(I) \subset \Phi_I$ est celle qu'on vient de donner. Il reste à prouver que si $I$ est noethérien, on a $\Phi_I \subset \Phi(I)$. Étant donnée $i : A \longrightarrow B$ dans $\Phi_I$, il faut prouver que $u \in \Phi(I)$, % [sic] u -> i, the arrow in question is i : A -> B
i.e. que pour un diagramme commutatif dans $\mathcal{M}(I)$,
\[
\begin{tikzcd}
A \arrow[r, "\alpha"] \arrow[d, "i"'] & X \arrow[d, "p"] \\
B \arrow[r, "\beta"'] \arrow[ur, dotted, "h"] & Y
\end{tikzcd}
\tag*{(6.3.2)}
\]
avec $p \in \Psi\langle I\rangle$, i.e. $p(i) : X(i) \longrightarrow Y(i)$ dans $\Psi$ pour tout $i$, $\exists\ h : B \longrightarrow X$ laissant le diagramme commutatif. On va revenir aux notations duales, en travaillant dans $\mathcal{M}(I^{\circ})$, avec $I$ artinien. On considère l'ensemble $\mathcal{H}$ des couple % [sic] des couple -> des couples, plural
$(U, h_U)$ d'un ouvert $U$ de $I$ et d'un $h_U : B|U \longrightarrow X|U$, rendant commutatif la restriction à $U$ de (6.3.2). % [sic] commutatif -> commutative, agreement with « la restriction »
On ordonne $\mathcal{H}$ par prolongement, on note que toute partie totalement ordonnée $(U_\alpha, h_{U_\alpha})$ dans $\mathcal{H}$ a une borne supérieure, savoir $U = \bigcup_\alpha U_\alpha$, $h_U$ défini par les $h_{U_\alpha}$ sur les $U_\alpha$.
% foot: Version 9 février 2019

\printedpage{55}\msp{page 110}
% head: A. Grothendieck : Dérivateurs, Ch. XVII, Édition des Universités de Montpellier II et Paris VII
(\textbf{N.B.} L'ensemble $\mathcal{H}$ n'est pas vide, car $(\emptyset, h_{\emptyset})$ est dans $\mathcal{H}$ -- mais c'est contenu dans l'énoncé précédent que toute partie totalement ordonnée de $\mathcal{H}$ y a une borne supérieure \dots). Par Zorn, il existe un élément maximal $(U, h_U)$. À prouver que $U = I$, et on aura terminé. Ou encore que si $(U, h_U)$ tel que $U \neq I$, ce n'est \textit{pas} maximal. Soit $F = I\backslash U \neq \emptyset$ le fermé complémentaire. Par l'hypothèse artinienne sur $I$, $F$ a un élément minimal $i_0$. Soit
\[
U' \ =\ U \cup \overbrace{I/i_0}^{V} \ =\ U \cup \{i_0\}\ ,
\]
on va prolonger $h_U$ en $h_{U'}$ sur $U'$, on aura terminé. Pour ceci, soit $V_0 = V \cap U = V\backslash\{i_0\}$, on est ramené à prolonger la restriction $h_{V_0}$ de $h_U$ à $V_0$, de $V_0$ à $V$. Au total, on peut supposer $I = I/i_0$ muni d'un élément final $i_0$, $U = I\backslash\{i_0\}$, et montrer qu'une $(U, h_U) \in \mathcal{H}$ se prolonge à $I$. Soient $A'$, $B'$, $X'$, $Y'$ les restrictions de $A$, $B$, $X$, $Y$ à $U$ (qu'on pourra aussi noter $I'$). La donné de $A$ % [sic] La donné -> La donnée, missing feminine ending
équivaut à celle de $A'$, celle de $A_0 = A(i_0) \in \mathrm{Ob}\,\mathcal{M}$, et d'une flèche de
\[
A'_{I'} \ \overset{\text{déf}}{=}\ \varinjlim\limits_{i' \in I'} A'(i') \ \xrightarrow{\ \lambda_A\ }\ A_0\ (\ \simeq A_I\ )\ ,
\]
\msp{page 111}et de même pour $B$, $X$, $Y$. Plus précisément, la catégorie $\mathcal{M}(I^{\circ})$ s'interprète comme celle des triples $X = (X', X_0, \lambda_X)$, où $X' \in \mathcal{M}(I'^{\circ})$, $X_0 \in \mathcal{M}$, et $\lambda_X : X'_{I'} \longrightarrow X_0$. À ce compte, la donnée du diagramme (6.3.2) équivaut à celle d'un diagramme carré similaire
\[
\begin{tikzcd}
A' \arrow[r, "\alpha'"] \arrow[d, "i'"'] & X' \arrow[d, "p'"] \\
B' \arrow[r, "\beta'"'] \arrow[ur, dotted, "h'"] & Y'
\end{tikzcd}
\tag*{(6.3.3)}
\]
\textit{dans} $\mathcal{M}(I'^{\circ})$, d'où un diagramme dans $\mathcal{M}$
\[
\begin{tikzcd}
A'_{I'} \arrow[r, "\alpha'_{I'}"] \arrow[d, "i'_{I'}"'] & X'_{I'} \arrow[d, "p'_{I'}"] \\
B'_{I'} \arrow[r, "\beta'_{I'}"'] \arrow[ur, dotted, "h'_{I'}"] & Y'_{I'}\ ,
\end{tikzcd}
\tag*{(6.3.4)}
\]
d'un autre diagramme carré \textit{dans} $\mathcal{M}$
\[
\begin{tikzcd}
A_0 \arrow[r, "\alpha_0"] \arrow[d, "i_0"'] & X_0 \arrow[d, "p_0"] \\
B_0 \arrow[r, "\beta_0"'] & Y_0\ ,
\end{tikzcd}
\tag*{(6.3.5)}
\]
et enfin d'un homomorphisme $\lambda = (\lambda_A, \lambda_B, \lambda_X, \lambda_Y)$ du carré (6.3.4) dans le carré (6.3.5) :
% foot: Version 9 février 2019
\printedpage{56}\msp{page 112}
% head: A. Grothendieck : Dérivateurs, Ch. XVII, Édition des Universités de Montpellier II et Paris VII
\[
\begin{tikzcd}
A'_{I'} \arrow[rr, "\alpha'_{I'}"] \arrow[dr, "i'_{I'}"'] \arrow[dd, "\lambda_A"'] & & X'_{I'} \arrow[dr, "p'_{I'}"] \arrow[dd, "\lambda_X" near start] & \\
 & B'_{I'} \arrow[ur, dotted, "h'_{I'}"] \arrow[rr, crossing over, "\beta'_{I'}" near end] & & Y'_{I'} \arrow[dd, "\lambda_Y"] \\
A_0 \arrow[rr, "\alpha_0" near start] \arrow[dr, "i_0"'] & & X_0 \arrow[dr, "p_0"] & \\
 & B_0 \arrow[from=2-2, crossing over, "\lambda_B"' near start] \arrow[ur, dotted, "h_0"'] \arrow[rr, "\beta_0"] & & Y_0
\end{tikzcd}
\qquad\qquad \text{\footnotemark[32]}
\tag*{(6.3.6)}
\]
\footnotetext[32]{\textbf{N.B.} La condition $p \in \Psi\langle I\rangle$ s'exprime par $p' \in \Psi\langle I'\rangle$ et \fbox{$p_0 \in \Psi.$}}

D'autre part, la question de prolonger $(I', h') \in \mathcal{H}$ en $(I, h)$, revient à celle de trouver $h_0 : B_0 \longrightarrow X_0$ dans $\mathcal{M}$, de façon qu'avec $h'_{I'}$ cela fasse encore un diagramme commutatif de (6.3.6), ce qui revient à dire que d'une part les deux triangles formés dans le carré de base sont commutatifs :
\[
p_0h_0 \ =\ \beta_0, \qquad\quad h_0i_0 \ =\ \alpha_0\ ,
\tag*{(6.3.7)}
\]
et d'autre part le carré diagonal défini par $h'_{I'}$, $h_0$ soit commutatif (ce qui exprime qu'on a une flèche $h$ de $B$ dans $X$) :
\[
\lambda_X h'_{X'} \ =\ h_0\lambda_B\ . % [sic] h'_{X'} -> h'_{I'}, the arrow is h'_{I'} everywhere else (cf. (6.3.16))
\tag*{(6.3.8)}
\]

\msp{page 113}Mais le diagramme (6.3.6), où la flèche en traits pointillés denses $h'_{I'}$ est déjà acquise, et $h_0$ en pointillés espacés reste à trouver, donne naissance au diagramme, dans lequel est mis en évidence la flèche $\tilde{i} \in \Phi$ qui exprime l'hypothèse $i \in C(i)$, % [sic] C(i) -> C(I), cf. i' \in C(I')
compte tenu de $i' \in C(I')$ :
\[
\begin{tikzcd}
A'_{I'} \arrow[rr, "\alpha'_{I'}"] \arrow[dr, "i'_{I'}"'] \arrow[dd, "\lambda_A"'] & & X'_{I'} \arrow[dr, "p'_{I'}"] \arrow[dd, "\lambda_X" near start] & \\
 & B'_{I'} \arrow[ur, dotted, "h'_{I'}"] \arrow[rr, crossing over, "\beta'_{I'}" near end] & & Y'_{I'} \arrow[dd, "\lambda_Y"] \\
A_0 \arrow[ur, phantom, "\text{cocart.}"] \arrow[rr, "\alpha_0" near start] \arrow[dr, "r"'] \arrow[ddr, dotted, "i_0"'] & & X_0 \arrow[dr, "p_0\in\Psi"] & \\
 & \tilde{B} \arrow[from=2-2, crossing over, "s"' near start] \arrow[ur, dotted, "\tilde{h}_0"] \arrow[rr, dotted, "\tilde{\beta}_0"] \arrow[d, "\tilde{i}\in\Phi"'] & & Y_0 \\
 & B_0\ , \arrow[uur, dotted, "h_0"'] \arrow[urr, "\beta_0"'] & &
\end{tikzcd}
\tag*{(6.3.9)}
\]
% foot: Version 9 février 2019
\printedpage{57}
% head: A. Grothendieck : Dérivateurs, Ch. XVII, Édition des Universités de Montpellier II et Paris VII
où la flèche $\tilde{h}_0$ en traits pointillés denses de $\tilde{B} \overset{\text{déf}}{=} A_0 \amalg_{A'_{I'}} B'_{I'}$ dans $X_0$, est définie par les conditions de commutativité
\[
\tilde{h}_0r \ =\ \alpha_0, \qquad\quad \tilde{h}_0s \ =\ \lambda_X h'_{I'}\ ,
\tag*{(6.3.10)}
\]
ce qui est licite, vu que
\[
\begin{aligned}
\alpha_0\lambda_A \ &=\ (\lambda_X h'_{I'})i'_{I'} \\
&=\ \lambda_X(h'_{I'}i'_{I'}) \\
&=\ \lambda_X\alpha'_{I'}.
\end{aligned}
\]

\msp{page 114}Quant à la flèche $\tilde{\beta}_0 : \tilde{B} \longrightarrow Y_0$, elle est donnée par la condition de commutativité
\[
\tilde{\beta}_0r \ =\ p_0\alpha_0, \qquad\quad \tilde{\beta}_0s \ =\ \lambda_Y\beta'_{I'}\ ,
\tag*{(6.3.11)}
\]
ce qui est encore licite, puisque
\[
(p_0\alpha_0)\lambda_A \ =\ (\lambda_Y\beta'_{I'})i'_{I'}
\]
en vertu de la commutativité du cube (6.3.6). Pour établir la commutativité totale du diagramme (6.3.9) (traits pleins, et traits pointillés denses pour $h'_{I'}$, $\tilde{h}_0$ et $\tilde{\beta}_0$ -- à l'exclusion de $h_0$ pas encore construit), il faut encore vérifier d'une part les relation % [sic] les relation -> les relations, plural
\[
\left\{
\begin{aligned}
\tilde{h}_0r \ &=\ \alpha_0, \quad p_0\tilde{h}_0 \ =\ \tilde{\beta}_0, \\
\lambda_X h'_{I'} \ &=\ \tilde{h}_0s
\end{aligned}
\right.
\qquad\qquad \text{\footnotemark[33]},
\tag*{(6.3.12)}
\]
\footnotetext[33]{\textbf{N.B.} la première et la troisième de ces relations sont déjà connues, cf. (6.3.10).}
établissant la commutativité du cube supérieur avec le carré diagonal (diagramme du type (6.3.6), la flèche pointillée $h_0$ y incluse), d'autre part les relations
\[
\left\{
\begin{aligned}
\tilde{i}r \ &=\ i_0 \qquad \text{(et $\tilde{i}s = \lambda_B$ pour mémoire)} \\
\beta_0\tilde{i} \ &=\ \tilde{\beta}_0\ ,
\end{aligned}
\right.
\tag*{(6.3.13)}
\]
où la première ligne est déjà connue par définition de $\tilde{B}$ et de $\tilde{i}$. La deuxième
\msp{page 115}se ramène aux deux relations
\[
\begin{aligned}
(\beta_0\tilde{i})r \ &=\ \beta_0(\tilde{i}r) \\
&=\ \beta_0i_0 \\
&=\ p_0\alpha_0 \\
&=\ \tilde{\beta}_0r\ ,
\end{aligned}
\]
OK par commutativité du carré de base de (6.3.5),
\[
\begin{aligned}
(\beta_0\tilde{i})s \ &=\ \beta_0(\tilde{i}s) \\
&=\ \beta_0\lambda_B \\
&=\ \lambda_Y\beta'_{I'} \\
&=\ \tilde{\beta}_0s\ ,
\end{aligned}
\]
% foot: Version 9 février 2019
\printedpage{58}
% head: A. Grothendieck : Dérivateurs, Ch. XVII, Édition des Universités de Montpellier II et Paris VII
OK par commutativité du carré antérieur dans (6.3.6).

Donc il reste à vérifier la deuxième relation (6.3.12),
\[
p_0\tilde{h}_0 \ =\ \tilde{\beta}_0\ ,
\]
ce qui se ramène à
\[
\begin{aligned}
(p_0\tilde{h}_0)r \ &=\ p_0(\tilde{h}_0r) \\
&=\ p_0\alpha_0 \\
&=\ \tilde{\beta}_0r\ ,
\end{aligned}
\]
OK par (6.3.11),
\[
\begin{aligned}
(p_0\tilde{h}_0)s \ &=\ p_0(\tilde{h}_0s) \\
&=\ p_0(\lambda_X h'_{I'}) \\
&=\ (p_0\lambda_X)h'_{I'} \\
&=\ (\lambda_Y p'_{I'})h'_{I'} \\
&=\ \lambda_Y(p'_{I'}h'_{I'}) \\
&=\ \lambda_Y\beta'_{I'} \\
&=\ \tilde{\beta}_0s\ ,
\end{aligned}
\]
OK.

Ceci posé, je dis que les contraintes (6.3.7), (6.3.8) sur $h_0$ reviennent à la commutativité des deux triangles dans les carré % [sic] les carré -> le carré, singular (one square (6.3.14))
\[
\begin{tikzcd}
\tilde{B} \arrow[r, "\tilde{h}_0"] \arrow[d, "\tilde{i}\in\Phi"'] & X_0 \arrow[d, "p_0\in\Psi"] \\
B_0 \arrow[r, "\beta_0"'] \arrow[ur, dotted, "h_0"] & Y_0\ ,
\end{tikzcd}
\tag*{(6.3.14)}
\]
\msp{page 116}i.e. aux deux relations
\[
h_0\tilde{i} \ =\ \tilde{h}_0\ , \qquad\quad p_0h_0 \ =\ \beta_0\ .
\tag*{(6.3.15)}
\]

La deuxième de ces relations n'est autre que la première relation (6.3.7), et il reste à voir que la première relation (6.3.15) équivaut à l'ensemble des deux relations restantes dans (6.3.7), (6.3.8) :
\[
h_0i_0 \ =\ \alpha_0\ , \qquad\quad \lambda_X h'_{I'} \ =\ h_0\underbrace{\lambda_B}_{\tilde{i}s}
\tag*{(6.3.16)}
\]
% foot: Version 9 février 2019

\end{document}
